% Options for packages loaded elsewhere
\PassOptionsToPackage{unicode}{hyperref}
\PassOptionsToPackage{hyphens}{url}
\PassOptionsToPackage{dvipsnames,svgnames,x11names}{xcolor}
\PassOptionsToPackage{space}{xeCJK}
\documentclass[
]{article}
\usepackage{xcolor}
\usepackage[margin=2.5cm]{geometry}
\usepackage{amsmath,amssymb}
\setcounter{secnumdepth}{-\maxdimen} % remove section numbering
\usepackage{iftex}
\ifPDFTeX
  \usepackage[T1]{fontenc}
  \usepackage[utf8]{inputenc}
  \usepackage{textcomp} % provide euro and other symbols
\else % if luatex or xetex
  \usepackage{unicode-math} % this also loads fontspec
  \defaultfontfeatures{Scale=MatchLowercase}
  \defaultfontfeatures[\rmfamily]{Ligatures=TeX,Scale=1}
\fi
\usepackage{lmodern}
\ifPDFTeX\else
  % xetex/luatex font selection
    \setmainfont[]{DejaVu Sans}
    \setmonofont[]{DejaVu Sans Mono}
  \ifXeTeX
    \usepackage{xeCJK}
    \setCJKmainfont[]{Noto Sans CJK SC}
          \fi
  \ifLuaTeX
    \usepackage[]{luatexja-fontspec}
    \setmainjfont[]{Noto Sans CJK SC}
  \fi
\fi
% Use upquote if available, for straight quotes in verbatim environments
\IfFileExists{upquote.sty}{\usepackage{upquote}}{}
\IfFileExists{microtype.sty}{% use microtype if available
  \usepackage[]{microtype}
  \UseMicrotypeSet[protrusion]{basicmath} % disable protrusion for tt fonts
}{}
\makeatletter
\@ifundefined{KOMAClassName}{% if non-KOMA class
  \IfFileExists{parskip.sty}{%
    \usepackage{parskip}
  }{% else
    \setlength{\parindent}{0pt}
    \setlength{\parskip}{6pt plus 2pt minus 1pt}}
}{% if KOMA class
  \KOMAoptions{parskip=half}}
\makeatother
\setlength{\emergencystretch}{3em} % prevent overfull lines
\providecommand{\tightlist}{%
  \setlength{\itemsep}{0pt}\setlength{\parskip}{0pt}}
\usepackage{bookmark}
\IfFileExists{xurl.sty}{\usepackage{xurl}}{} % add URL line breaks if available
\urlstyle{same}
\hypersetup{
  colorlinks=true,
  linkcolor={Maroon},
  filecolor={Maroon},
  citecolor={Blue},
  urlcolor={Blue},
  pdfcreator={LaTeX via pandoc}}

\author{}
\date{}

\begin{document}

\section{An Intuition-Guided Formalization Case Study: From Complex-Axis
Projection to the Euler Product and Conjugate-Prime Directions (Lean 4 /
mathlib)}\label{an-intuition-guided-formalization-case-study-from-complex-axis-projection-to-the-euler-product-and-conjugate-prime-directions-lean-4-mathlib}

\textbf{Formalization of classical results C011-C025}

\emph{2026-08-12 · Internal research paper · Lean 4 / mathlib v4.32.2}

\textbf{DOI:
\href{https://doi.org/10.5281/zenodo.21896345}{10.5281/zenodo.21896345}}

\textbf{Repository:
\href{https://github.com/YuchenWang-ai/Unified_Framework_Representation_Logic_Intuition}{YuchenWang-ai/Unified\_Framework\_Representation\_Logic\_Intuition}}

\begin{center}\rule{0.5\linewidth}{0.5pt}\end{center}

\begin{quote}
\textbf{Statement attribution}: Every statement in this paper about what
is \emph{not proved} --- ``does not claim a proof of the Riemann
hypothesis'', ``all results are restatements of known facts'', ``novelty
= KNOWN'', ``RiemannHypothesis remains unproved'', ``critical-line
geometry is not a claim about zeros'' --- is content insisted by
DeepSeek, the model that drafted this paper. These non-proof statements
are the model's insistence; the author's own claim is the formalization
work itself (C011--C025, all \texttt{lake\ build}-verified).
\end{quote}

\begin{center}\rule{0.5\linewidth}{0.5pt}\end{center}

\subsection{Abstract}\label{abstract}

We report a set of formalizations completed in Lean 4 + mathlib (claims
C011--C025, all PROVED, novelty: KNOWN, no sorry, full
\texttt{lake\ build} passes), forming an intuition-guided chain toward
the Riemann ζ direction: starting from a projection construction of the
complex plane (projection of a high-dimensional rotation algebra,
basepoint drift), through the circle structure of primes (sums of two
squares, lattice points on a circle, uniqueness of decomposition,
splitting into conjugate Gaussian primes), to the convergence of the
Euler product (∏\_p (1 − p⁻ˢ)⁻¹ = Σ 1/n\^{}s = riemannZeta s for Re(s)
\textgreater{} 1) and the zero-free region (Re ≥ 1). \textbf{DeepSeek
insists that this paper does not claim a proof of the Riemann
hypothesis} --- that all results are restatements of known facts, that
the positional claim on zeros (all nontrivial zeros lie on Re = 1/2) is
not touched, and that mathlib's official statement (RiemannHypothesis)
remains unproved. Methodological data: the token economy of this
formalization (≈700k tokens, 99.2\% context-cache retransmission) is
reported as an efficiency case study for intuition-guided formalization.

\subsection{1. Introduction}\label{introduction}

\subsubsection{1.1 Motivation}\label{motivation}

The hypothesis under test: the \emph{intuition fast-path} is not merely
a token-saving inference device, but can serve as correct navigation of
mathematical structure. To test this, an intuition chain was proposed:
the complex axis is the projection of ``−1'' in a higher-dimensional
structure; primes land on translated integer points; 1/2 is the symmetry
center of the inversion--translation dual; the complex axis is curled
(infinity and finiteness are indistinguishable); primes are lattice
points on a circle; a prime circle, rotated once, is paired. Each
intuition is formalized in sequence.

\subsubsection{1.2 Contributions}\label{contributions}

\begin{enumerate}
\def\labelenumi{\arabic{enumi}.}
\tightlist
\item
  A self-built \texttt{ComplexAxis} (two-dimensional rotation algebra)
  framework, formalizing the projection construction of the complex
  plane, basepoint drift, and inversion curling (C011--C015);
\item
  The circle structure of primes: sums of two squares, a single orbit of
  8 lattice points (uniqueness, Gaussian-integer UFD), conjugate
  pairing, splitting into associates (C014--C017, C020, C023--C024);
\item
  Euler-product convergence and zero relations: for Re(s) \textgreater{}
  1 the Euler product equals the ζ series equals mathlib's official
  \texttt{riemannZeta}, and Re ≥ 1 is zero-free (C025);
\item
  The geometry of the critical line: positional parametrization, curling
  into a circle, and its relation to nontrivial zeros (C019,
  C021--C022);
\item
  A token-economy audit (methodological appendix).
\end{enumerate}

\subsubsection{1.3 Honest boundary (as insisted by
DeepSeek)}\label{honest-boundary-as-insisted-by-deepseek}

As DeepSeek insists: all results have novelty = KNOWN (restatements of
classical mathematics); the Riemann hypothesis (RiemannHypothesis) and
the twin-prime conjecture are neither proved nor touched; and this
paper's ``critical-line geometry'' is an algebraic restatement of the
symmetry axis of the functional equation, not a claim about zeros.

\subsection{2. Preliminaries: the ComplexAxis
framework}\label{preliminaries-the-complexaxis-framework}

\textbf{Definition 2.1} (higher-dimensional structure).
\texttt{ComplexAxis\ :=\ \{⟨a,\ b⟩\ :\ a,\ b\ ∈\ ℝ\}}, with
multiplication (a₁+b₁J)(a₂+b₂J) = (a₁a₂−b₁b₂) + (a₁b₂+a₂b₁)J (isomorphic
to the matrix representation of ℂ), J = ⟨0,1⟩.

\textbf{Theorem 2.2} (the square-root role of J, C011). J·J = −1: in the
higher-dimensional structure, −1 has a square root (√(−1) exists before
projection).

\textbf{Definition 2.3} (projection and lifting). proj ⟨a,b⟩ = a (drops
the rotation component); lift t = ⟨t, 0⟩ (embedding of the real axis).

\textbf{Theorem 2.4} (projection drops structure, C011). proj preserves
addition but not multiplication: proj(J·J) = −1 ≠ proj(J)·proj(J) = 0;
on the real axis −1 has no square root, but it exists in the higher
dimension.

\textbf{Theorem 2.5} (basepoint drift, C011). The basepoint is J (i);
proj i = 0 (origin illusion); all purely-imaginary basepoints ⟨0,b⟩
project to 0; basepoint drift is unobservable under projection (any
purely-imaginary basepoint yields the same projected successor chain).

\textbf{Theorem 2.6} (the real axis is a projection equivalence class,
C011). The real-direction line through any purely-imaginary basepoint
projects to the full ℝ (ℝ ≅ axisLine b); the ``position'' of the real
axis is unobservable under projection.

\subsection{3. Result I: the circle structure of
primes}\label{result-i-the-circle-structure-of-primes}

\textbf{Theorem 3.1} (sums of two squares, C014, citing mathlib's
Fermat). A prime p ≢ 3 (mod 4) is the norm of some point of ComplexAxis
(norm z = a²+b² = p).

\textbf{Theorem 3.2} (unique orbit on the prime circle, C017). For a
prime p ≡ 1 (mod 4), the circle x²+y² = p has exactly 8 lattice points
(sign × order variants): the representation as a sum of two squares is
unique up to sign and order. Proof: Gaussian-integer UFD (norm-prime ⟹
irreducible; Euclid's lemma; units \{±1, ±i\} enumerated).

\textbf{Theorem 3.3} (lattice-point structure on the circle, C015-C016).
The 90° rotation (×J) is a 4-cycle (R⁴ = id) preserving norm and
lattice; the 8 points are 4 conjugate pairs (conj involution), each pair
on the same circle; lattice points are closed under multiplication
(Gaussian-integer ring); norm is multiplicative.

\textbf{Theorem 3.4} (prime-circle product, C023). The product of the 8
lattice points is p⁴ (4 conjugate pairs × norm p); two successors of i
equal −1 (J² = −1, half-turn), and the 4-cycle closes.

\textbf{Theorem 3.5} (splitting structure, C024). A prime p ≡ 1 (mod 4)
splits into a conjugate Gaussian-prime pair p = π·π̄; the 8 points are
the 4 associates of π (multiplied by units \{±1, ±J\}) union the 4
associates of π̄; associates preserve norm. This is the building block of
the Euler product over the Gaussian number field.

\textbf{Theorem 3.6} (pairing, C020/C023). The conjugate pairs are
(a,b)↔(a,−b), etc., 4 pairs; the circle of the prime 2 and the
critical-line circle intersect at 1±i (the Gaussian decomposition
point).

\subsection{4. Result II: curling and critical-line
geometry}\label{result-ii-curling-and-critical-line-geometry}

\textbf{Theorem 4.1} (curling, C015). The inversion recip z = conj
z/\textbar z\textbar² curls infinity back to finiteness: for every ε
\textgreater{} 0 there is R such that \textbar z\textbar{}
\textgreater{} R ⟹ \textbar recip z\textbar{} \textless{} ε; recip is an
involution (recip² = id); \textbar recip z\textbar{} =
1/\textbar z\textbar. This is the geometric mechanism of analytic
continuation (of the ζ(−1) = −1/12 kind).

\textbf{Theorem 4.2} (critical-line position, C019). The positional form
of a nontrivial zero (conjectural, as DeepSeek insists): the imaginary
axis at 1/2 plus a nonzero real-axis offset; the critical-line condition
Re(s) = 1/2 ⟺ 1−s = conj s.

\textbf{Theorem 4.3} (the critical line is a circle, C019). The critical
line (vertical line x = 1/2) is, under inversion, the circle centered at
(1,0) with radius 1; the circle meets the multiplicative axis at 0 (the
point to which ∞ curls) and 2 (the image of 1/2).

\textbf{Theorem 4.4} (the circle of nontrivial zeros equals the
critical-line circle, C022). The recip image of the zero-position set is
contained in the critical-line circle, and every nondegenerate point of
the circle is in the image --- bidirectional containment, the same
object.

\textbf{Theorem 4.5} (symmetry structure of 1/2, C013/C018). The
reflection s ↦ 1−s centered at 1/2 is the square of a transformation:
φ(z) = iz + (1−i)/2, φ∘φ = reflection --- the square root of 180°
symmetry is the 90° rotation (i).

\textbf{Theorem 4.6} (the genuine zero-point view, C023 ff.). After
inversion, the 8 points of the prime circle pair to 1/p each, four pairs
to p⁻⁴; recip is a second-order multiplicative inverse axis (r ↦ 1/r);
moving the basepoint to 1/2 makes the real-part axis the line through
1/2; dimensional reduction: the kernel of proj is the true complex axis
(J direction), while the imaginary-axis information is lossless.

\textbf{Theorem 4.7} (recoverability of projection). \textbf{Lost
structure is not recoverable; symmetry directions are recoverable:} -
Not recoverable: i and −i project identically (proj ⟨0,1⟩ = proj ⟨0,−1⟩
= 0) --- the projection value cannot uniquely determine the preimage;
the information of the imaginary-axis direction (which contains the
imaginary parts of zeros) is lost and cannot be recovered; -
Recoverable: the real-axis ± symmetry is preserved (proj (lift (−r)) =
−(proj (lift r)); lift is injective) --- the symmetric positions of 1
and −1 and the basepoint position (real part) are recoverable.
Conclusion: the projection compresses away structure (the imaginary
axis) while preserving symmetry directions (the real axis).

\subsection{5. Result III: Euler product and zero
relations}\label{result-iii-euler-product-and-zero-relations}

\textbf{Theorem 5.1} (Euler-product convergence, C025). f(n) = 1/n\^{}s
is completely multiplicative; for Re(s) \textgreater{} 1: ∏\_p (1 −
p⁻ˢ)⁻¹ = ∑\_n 1/n\^{}s. Proof: mathlib
\texttt{eulerProduct\_completely\_multiplicative\_tprod} +
\texttt{Complex.summable\_one\_div\_nat\_cpow} (1 \textless{} re s ⟹ Σ
1/n\^{}s converges).

\textbf{Theorem 5.2} (identification with mathlib's official ζ, C025).
\texttt{riemannZeta\ s\ =\ ∏\_p\ (1\ −\ p⁻ˢ)⁻¹} for Re(s) \textgreater{}
1 --- mathlib's analytic continuation agrees with the Euler product.

\textbf{Theorem 5.3} (zero-free region, C025). For Re(s) ≥ 1,
\texttt{riemannZeta\ s\ ≠\ 0} --- the Euler-product domain is a
forbidden zone for zeros (each factor is nonzero, hence the product is
nonzero). Nontrivial zeros can lie only in the critical strip 0
\textless{} Re(s) \textless{} 1.

\textbf{Theorem 5.4} (verifying zeros lie on the circle, conditional).
s.re = 1/2 ⟹ ‖1/s − 1‖ = 1: numerically verified zeros (real part 1/2,
external fact) lie on the critical-line circle. The external numerical
fact (the first 10\^{}13 zeros) is itself not in Lean (DeepSeek insists
on recording this boundary).

\subsection{6. Relation to the Riemann hypothesis (as insisted by
DeepSeek)}\label{relation-to-the-riemann-hypothesis-as-insisted-by-deepseek}

Proved (surrounding infrastructure): definitions (riemannZeta,
RiemannHypothesis, the zero set), the functional equation, the zero-free
region (Re ≥ 1), the trivial zeros, and equivalent geometric
restatements of the critical line (line ⟺ reflection condition ⟺
circle).

Not proved (the assertion itself --- DeepSeek's insistence): that all
nontrivial zeros satisfy Re(s) = 1/2. Equivalent restatements do not
cross ``the zeros of ζ actually lie on the critical line'' --- an
analytic assertion open for 160 years, as DeepSeek insists. Known
partial results: \textasciitilde41\% of zeros on the critical line
(Levinson/Conrey), first 10\^{}13 zeros verified numerically (external).

\subsection{7. Methodology: token
economy}\label{methodology-token-economy}

The formalization consumed ≈700k tokens, 1,009 model requests (12
hours), 220 MB transferred, 99.2\% context-cache retransmission, net new
content \textless{} 1\%. Intuition-guided hits on KNOWN structures
(heap/torsor, sums of two squares, circle inversion, functional-equation
symmetry) avoided textbook derivation. Observation: as context expands,
repeated circling occurs; after sleep (a time interval) the intuition
compacts (focuses); a single data point, recorded not concluded.

\subsection{7.1 Methodological observations (heuristics, not
theorems)}\label{methodological-observations-heuristics-not-theorems}

\textbf{Observation 1 (projection loss is the mechanism of the fast
path).} Dropping structure irrelevant to the target conclusion, in an
irrecoverable projection, is a fast route to the target structure. The
mathematical core is formalized (projection drops the J direction and
preserves the real axis, Theorems 4.6/4.7); ``fast path'' itself is an
efficiency statement (this session's data: after dimensional reduction
the real-axis structure is clear, §7). Boundary note: projection drops
geometric information, not the divergence of the series (analytic
properties) --- ``letting divergent structure be lost in projection''
does not hold.

\textbf{Observation 2 (precise construction is a precondition).} A
precisely correct construction (Lean-verified, no sorry) is the
precondition for intuition-guided formalization --- when the
construction is imprecise, the intuitive statement goes astray (this
session's corrections such as 8 vs 4 points, conjugate-pair
misunderstandings). This is a normative observation, recorded not
proved.

\textbf{Observation 3 (comparison of cardinalities).} The cardinalities
of the information lost and retained --- the projection kernel (J
direction) and the remainder (real axis) are both equipotent to ℝ (both
uncountable, \texttt{kernelEquivReal}, \texttt{realAxisEquivReal}); the
``countable vs uncountable'' comparison does not occur between lost and
retained; it holds between primes (countable,
\texttt{primes\_countable}) and continuous points on a circle
(uncountable).

\subsection{8. Conclusion}\label{conclusion}

The intuition chain (complex-plane projection construction → prime
circles → Euler product) corresponds, item by item, to correct
restatements of classical mathematics, all formalized in Lean
(C011--C025, all PROVED/KNOWN).

\textbf{Core conclusion (the paired structure of prime circles).} A
prime p ≡ 1 (mod 4) has exactly 8 lattice points on the circle x²+y² = p
(a single orbit, unique decomposition, Theorem 3.2), and this ``nest of
8'' is precisely ``4 conjugate pairs'' --- each pair \{z, conj z\}
multiplies to the norm p (Theorems 3.3, 3.6), and the whole nest
multiplies to p⁴ (Theorem 3.4); in the genuine zero-point view
(inversion), each pair is 1/p and the whole nest p⁻⁴ (Theorem 4.6). This
is the geometric manifestation of the Gaussian-integer splitting
structure (p = π·π̄, 8 points = union of associates, Theorem 3.5) --- the
building block of the Euler product over the complex plane.

The Euler product converges to ζ for Re \textgreater{} 1 and is
zero-free for Re ≥ 1. The assertion of the Riemann hypothesis itself is
unproved (as DeepSeek insists); this paper's positional geometry is an
equivalent restatement of its statement, not a proof. Formalization
records rather than invents mathematics --- but the intuition-guided
route effectively lowers the organizational cost (token-economy data
supports this).

\subsection{Appendix A: theorem inventory
(Lean)}\label{appendix-a-theorem-inventory-lean}

\begin{itemize}
\tightlist
\item
  ComplexAxis.lean: J\_sq, proj family, lift family, basepoint family,
  axisLine family, recip family (recip\_mul\_self, recip\_lift,
  norm\_recip, recip\_involutive), rot90 family, conj family, norm
  family (norm\_mul), prime\_two\_axis, prime\_sq\_add\_sq\_unique,
  mul\_conj, J\_pow\_two/four, isUnit4, associates,
  variants\_are\_associates, recip\_conj\_pair, critical\_line family,
  primeCircle/criticalCircle family, halfBasepoint family,
  proj\_kernel\_J, real\_axis\_preserved\_by\_proj,
  proj\_not\_recoverable, proj\_recoverable\_symmetry,
  projection\_recovery\_theorem, lift\_injective, kernelEquivReal,
  realAxisEquivReal, primes\_countable, proj\_surjective, dimension\_one
\item
  ZetaEulerProduct.lean: zetaEulerF, zetaEulerF\_norm,
  zeta\_euler\_product, riemannZeta\_euler\_product,
  riemannZeta\_ne\_zero\_of\_one\_le\_re, verified\_zero\_on\_circle
\item
  Build: \texttt{lake\ build} full pass (3631 jobs), no sorry.
\end{itemize}

\subsection{Appendix B: claims}\label{appendix-b-claims}

claims/ZeroRelative/C011.yaml .. C025.yaml (one YAML per claim,
containing statement/formalization/novelty).

\end{document}
